\chapter{常用不等式与估计方法索引}
\label{chap:appendix-estimates}

\begin{chapterguide}
本附录集中列出全书反复使用的估计工具。每条公式同时注明对象范围、等号条件或常数依赖；使用时仍须核对相应和、积分与范数是否有限。
\end{chapterguide}

\section{三角不等式、反三角不等式与范数估计}

绝对值满足
\[
 |x+y|\le|x|+|y|,\qquad
 \bigl||x|-|y|\bigr|\le|x-y|.
\]
有限和版本为
\[
 \left|\sum_{k=1}^n x_k\right|\le\sum_{k=1}^n|x_k|.
\]
范数 \(\|\cdot\|\) 的定义包含正定性、绝对齐次性与三角不等式，所以
\[
 \bigl|\|x\|-\|y\|\bigr|\le\|x-y\|.
\]
在内积空间中，Cauchy--Schwarz 不等式
\[
 |\langle x,y\rangle|\le\|x\|\,\|y\|
\]
由 \(0\le\|x-ty\|^2\) 的二次判别式得到；等号当且仅当 \(x,y\) 线性相关。

估计和时，应先判断是否需要保留抵消。直接取绝对值稳健但可能损失阶数；Fourier 分析和振荡积分中，过早使用三角不等式会抹去关键抵消。

\section{Young、Hölder、Minkowski 与 Jensen 不等式}

若 \(p,q>1\) 且 \(1/p+1/q=1\)，则对 \(a,b\ge0\)，
\[
 ab\le\frac{a^p}{p}+\frac{b^q}{q}.
\]
这称为 Young 不等式，等号当且仅当 \(a^p=b^q\)。它可由凸函数 \(t\mapsto e^t\) 或加权算术--几何平均证明。

\begin{theorem}[Hölder 不等式]
对有限序列 \(x_k,y_k\) 与共轭指数 \(p,q>1\)，
\[
 \sum_k|x_ky_k|
 \le\left(\sum_k|x_k|^p\right)^{1/p}
    \left(\sum_k|y_k|^q\right)^{1/q}.
\]
在测度空间上，若右端有限，则
\[
 \int|fg|
 \le\left(\int|f|^p\right)^{1/p}
    \left(\int|g|^q\right)^{1/q}.
\]
\end{theorem}

\begin{proof}
若任一范数为零结论直接成立。否则把 \(x,y\) 分别除以其 \(p,q\) 范数，对每一项应用 Young 不等式后求和或积分。
\end{proof}

\begin{theorem}[Minkowski 不等式]
对 \(1\le p<\infty\)，
\[
 \|x+y\|_p\le\|x\|_p+\|y\|_p,
\qquad
 \|f+g\|_{L^p}\le\|f\|_{L^p}+\|g\|_{L^p}.
\]
\end{theorem}

\begin{proof}
\(p=1\) 是三角不等式。若 \(p>1\)，
\[
 \sum|x+y|^p\le
 \sum |x|\,|x+y|^{p-1}
 +\sum |y|\,|x+y|^{p-1}.
\]
对两项应用 Hölder，并除以 \(\|x+y\|_p^{p-1}\)。积分版本相同。
\end{proof}

若 \(\varphi\) 在凸集上凸，\(\lambda_i\ge0\)、\(\sum_i\lambda_i=1\)，则 Jensen 不等式为
\[
 \varphi\left(\sum_i\lambda_ix_i\right)
 \le\sum_i\lambda_i\varphi(x_i).
\]
严格凸时，所有正权对应的 \(x_i\) 相等才取等。概率与积分版本要求权重总质量为 \(1\) 且相关积分存在。

\section{离散和连续 Grönwall 不等式}

\begin{theorem}[离散 Grönwall 不等式]
设 \(u_n,b_n\ge0\)，且
\[
 u_n\le A+\sum_{k=0}^{n-1}b_ku_k,\qquad A\ge0.
\]
则
\[
 u_n\le A\prod_{k=0}^{n-1}(1+b_k)
 \le A\exp\left(\sum_{k=0}^{n-1}b_k\right).
\]
\end{theorem}

\begin{proof}
令 \(v_n=A+\sum_{k<n}b_ku_k\)。则 \(u_n\le v_n\)，且
\[
 v_{n+1}=v_n+b_nu_n\le(1+b_n)v_n.
\]
归纳得到乘积界，再用 \(1+t\le e^t\)。
\end{proof}

\begin{theorem}[连续 Grönwall 不等式]
设 \(u,b\) 在 \([0,T]\) 上非负可积，且
\[
 u(t)\le A+\int_0^t b(s)u(s)\,ds.
\]
则几乎处处
\[
 u(t)\le A\exp\left(\int_0^t b(s)\,ds\right).
\]
\end{theorem}

\begin{proof}
若 \(A>0\)，令 \(v(t)=A+\int_0^tb(s)u(s)\,ds\)。则 \(u\le v\) 且几乎处处 \(v'\le bv\)，所以
\[
 (\log v)'\le b.
\]
积分后得到结论。\(A=0\) 可先以 \(A+\varepsilon\) 代替 \(A\)，再令 \(\varepsilon\downarrow0\)。
\end{proof}

\section{截断、分层、分解与 \texorpdfstring{\(\varepsilon/3\)}{epsilon/3} 方法}

无界对象常先截断：
\[
 T_M(x)=\max\{-M,\min\{x,M\}\}.
\]
截断满足 \(|T_M(x)-T_M(y)|\le|x-y|\) 且 \(|T_M(x)|\le M\)。证明可先在有界对象上完成，再控制尾部 \(\{|x|>M\}\)。

分层估计把对象按大小拆成
\[
 \{2^k<|f|\le2^{k+1}\},
\]
把连续幅度问题化为几何级数。常见分解还包括“局部部分加远尾”“低频加高频”“好集加坏集”；分解阈值须与目标误差或尺度协调。

\(\varepsilon/3\) 方法处理三项误差：
\[
 \|x_n-x\|\le\|x_n-y_n\|+\|y_n-y\|+\|y-x\|.
\]
先固定辅助对象 \(y\) 控制第三项，再对该固定选择控制第二项，最后选择 \(n\) 控制第一项。数字 \(3\) 没有特殊地位；关键是选择顺序与常数总和小于目标误差。

\section{常数依赖关系和定量估计记法}

记号 \(A\lesssim B\) 表示存在常数 \(C>0\) 使 \(A\le CB\)。必须从上下文说明 \(C\) 允许依赖哪些参数；必要时写作
\[
 A\le C(p,d,\Omega)B.
\]
\(A\asymp B\) 表示双向控制。渐近记号
\[
 f=O(g),\qquad f=o(g)
\]
分别表示 \(|f/g|\) 在指定极限过程中有界、趋于零；必须写明极限变量与参数是否一致。

吸收法把
\[
 X\le A+\theta X,\qquad 0\le\theta<1
\]
改写为 \(X\le A/(1-\theta)\)。若 \(\theta\) 依赖参数，就要保证它在所研究范围内统一小于 \(1\)。

\section{Cesàro 平均与 Stolz--Cesàro 定理}

\begin{theorem}[Cesàro 平均]
若 \(a_n\to L\)，则
\[
 \frac1n\sum_{k=1}^n a_k\longrightarrow L.
\]
\end{theorem}

\begin{proof}
给定 \(\varepsilon>0\)，取 \(N\) 使 \(k\ge N\) 时 \(|a_k-L|<\varepsilon/2\)。固定前缀和除以 \(n\) 后趋于零，故充分大时
\[
 \left|\frac1n\sum_{k=1}^n(a_k-L)\right|
 \le\frac1n\sum_{k<N}|a_k-L|+\frac{\varepsilon}{2}<\varepsilon.
\]
\end{proof}

逆命题不成立，例如 \(a_n=(-1)^n\) 的 Cesàro 平均趋于 \(0\)，原列不收敛。

\begin{theorem}[Stolz--Cesàro]
设 \(b_n\) 严格递增且 \(b_n\to+\infty\)。若
\[
 \frac{a_{n+1}-a_n}{b_{n+1}-b_n}\longrightarrow L\in\R,
\]
则 \(a_n/b_n\to L\)。
\end{theorem}

\begin{proof}
给定 \(\varepsilon>0\)，从某项起
\[
 (L-\varepsilon)(b_{k+1}-b_k)
 \le a_{k+1}-a_k
 \le(L+\varepsilon)(b_{k+1}-b_k).
\]
从固定 \(N\) 求和到 \(n-1\)，再除以 \(b_n\)。固定的 \(a_N,b_N\) 项除以 \(b_n\) 后趋零，夹逼得到结论。
\end{proof}

\section*{习题}
\addcontentsline{toc}{section}{习题}

\begin{enumerate}[label=\textbf{C.\arabic*.}]
 \exerciseitem{C.1} 证明有限和三角不等式与范数的反三角不等式。
 \exerciseitem{C.2} 用二次多项式非负证明 Cauchy--Schwarz 不等式及等号条件。
 \exerciseitem{C.3} 从 Young 不等式推导有限序列 Hölder 不等式。
 \exerciseitem{C.4} 补全 Minkowski 不等式证明中指数 \((p-1)q=p\) 的计算与零范数情形。
 \exerciseitem{C.5} 证明有限 Jensen 不等式，并分析严格凸等号条件。
 \exerciseitem{C.6} 证明离散 Grönwall 不等式；再处理 \(u_n\le A_n+\sum_{k<n}b_ku_k\) 且 \(A_n\) 单调的版本。
 \exerciseitem{C.7} 严格处理连续 Grönwall 中 \(A=0\) 的极限步骤。
 \exerciseitem{C.8} 证明截断映射 \(T_M\) 是 \(1\)-Lipschitz。
 \exerciseitem{C.9} 为三项误差分别分配 \(\varepsilon/2,\varepsilon/4,\varepsilon/4\)，说明与 \(\varepsilon/3\) 的等价作用。
 \exerciseitem{C.10} 证明 Cesàro 平均定理，并构造逆命题反例。
 \projectexerciseitem{C.11} \ 【前置：单调性与求和分部；预计用时：60分钟】证明 Stolz--Cesàro 定理的有限极限及 \(+\infty\) 版本。
 \optionalexerciseitem{C.12} 对调和数 \(H_n\) 应用 Stolz--Cesàro，证明 \(H_n/\log n\to1\)；可使用 \(\log(1+1/n)\sim1/n\)。
\end{enumerate}
